\documentclass[10pt]{exam}
\printanswers
\usepackage{epsfig,amssymb}
\usepackage{xcolor}
\definecolor{darkred}{rgb}{0.5,0,0}
\definecolor{darkgreen}{rgb}{0,0.5,0}
\usepackage{hyperref}
\usepackage{fullpage}
\usepackage{tikz}
\pagestyle{empty} 
\usepackage{listings}
\setlength{\parindent}{0pt} %
\setlength{\parskip}{.25cm}
\newcommand{\comment}[1]{}
\boxedpoints
\addpoints
\usepackage{amsmath}
\usepackage{algorithm2e}
\usepackage{url}
\usepackage{enumitem}
\usepackage{graphics}

\pagestyle{headandfoot}
\firstpageheader{\Large CSCE 235}{{\Large Assignment 5}}{\Large Spring 2026}

\begin{document}

\hrule

\vspace{.50cm} Name:
\text{Charlie Pyle} 
\\ NUID: 22821006
\vspace{.50cm}

\textbf{Instructions} Follow instructions \emph{carefully}; failure to do so will result in points being deducted. 

\begin{questions}

\question[16] Using the \textbf{formal definition of Big-O}, show that each of these functions is $O(x^2)$. You must clearly state what the witnesses to the relationship are to receive full credit.
\begin{parts}
  \part $f(x)=26x$
  \part $f(x)=20x+26$
  \part $f(x)=20x^2+2x+6$
  \part $f(x)=20x^2+26log(x)$
\end{parts}

\begin{solution}
    Enter your solution here:\\
    \textbf{A:}\\
    For $x \ge 1$ we have $x \le x^2$, so
    \[26x \;\le\; 26x^2.\]
    Thus $f(x) \in O(x^2)$ with witnesses $c = 26$, $k = 1$.
 
    \textbf{B:}\\
    For $x \ge 1$ we have $x \le x^2$ and $1 \le x^2$, so
    \[20x + 26 \;\le\; 20x^2 + 26x^2 \;=\; 46x^2.\]
    Thus $f(x) \in O(x^2)$ with witnesses $c = 46$, $k = 1$.
 
    \textbf{C:}\\
    For $x \ge 1$ we have $x \le x^2$ and $1 \le x^2$, so
    \[20x^2 + 2x + 6\;\le\; 20x^2 + 2x^2 + 6x^2\;=\; 28x^2.\]
    Thus $f(x) \in O(x^2)$ with witnesses $c = 28$, $k = 1$.
 
    \textbf{D:}\\
    For $x \ge 1$ we have $\log(x) \le x \le x^2$, so
        \[20x^2 + 26\log(x)\;\le\; 20x^2 + 26x^2\;=\; 46x^2.\]
    Thus $f(x) \in O(x^2)$ with witnesses $c = 46$, $k = 1$.
\end{solution}

\question[8] Using the \textbf{formal definition of Big-O}, show that $20xlogx$ is $O(26x^2)$ but that $26x^2$ is not $O(20xlog(x))$. You must clearly state what the witnesses to the relationship are to receive full credit.

\begin{solution}
    Enter your solution here:\\
    \textbf{Part 1: $20x\log x \in O(26x^2)$}
    For $x \ge 1$ we have $\log(x) \le x$, therefore
    \[20x\log x \;\le\; 20x \cdot x \;=\; 20x^2 \;\le\; 26x^2.\]
    Thus $20x\log x \in O(26x^2)$ with witnesses $c = 1$, $k = 1$.
 
    \textbf{Part 2: $26x^2 \notin O(20x\log x)$}
    Suppose for contradiction that $26x^2 \in O(20x\log x)$. Then there exist constants $c > 0$ and $k$ such that $26x^2 \le c \cdot 20x\log x$ for all $x \ge k$, which simplifies to
    \[\frac{26x}{20c\log x} \;\le\; 1 \quad\text{for all } x \ge k.\]
    However,
    \[\lim_{x \to \infty} \frac{26x}{20c\log x}
        \;=\; \frac{1}{c}\lim_{x \to \infty} \frac{26x}{20\log x}\;=\; \infty,\]
    a contradiction. Thus $26x^2 \notin O(20x\log x)$.
\end{solution}

\question[16]  For each of the following functions, give a proof using the \textbf{Limit method} that $f(x) \in O(g(x))$:
\begin{parts}
\part $(f(x) = 20x^2, g(x) = 26x^3 $
\part $f(x) = 20xlog(x), g(x) = 26x^2$
\part $f(x) = 20 + log(x) + 2x^2 + 6x^3, g(x) = x^3$
\part $f(x) = 20xlog(x) + 2x, g(x) = 6x^2$
\end{parts}

\begin{solution}\\
    Enter your solution here:\\
    \textbf{A:}
    \[\lim_{x \to \infty} \frac{20x^2}{26x^3}
        \;=\; \frac{20}{26}\lim_{x \to \infty} \frac{1}{x}\;=\; 0.\]
    Since the limit is $0 < \infty$, we conclude $f(x) \in O(g(x))$.
 
    \textbf{B:}
    Divide numerator and denominator by $x$ (valid for $x > 0$):
    \[
        \lim_{x \to \infty} \frac{20x\log x}{26x^2}
        \;=\; \lim_{x \to \infty} \frac{20\log x}{26x}.
    \]
    Apply L'H\^{o}pital's Rule:
    \[\lim_{x \to \infty} \frac{20\log x}{26x}
        \;=\; \lim_{x \to \infty} \frac{\tfrac{20}{x}}{26}
        \;=\; \lim_{x \to \infty} \frac{20}{26x}\;=\; 0.\]
    Since the limit is $0 < \infty$, we conclude $f(x) \in O(g(x))$.
 
    \textbf{C:}
    \[\lim_{x \to \infty} \frac{20 + \log x + 2x^2 + 6x^3}{x^3}
        \;=\;\lim_{x \to \infty}\frac{20}{x^3}+ \lim_{x \to \infty}\frac{\log x}{x^3}+ \lim_{x \to \infty}\frac{2}{x}
        + \lim_{x \to \infty} 6\;=\; 0 + 0 + 0 + 6\;=\; 6.\]
    Since the limit is $6 < \infty$, we conclude $f(x) \in O(g(x))$.
 
    \textbf{D:}
    \[\lim_{x \to \infty} \frac{20x\log x + 2x}{6x^2}
        \;=\;\lim_{x \to \infty} \frac{20\log x}{6x}
        + \lim_{x \to \infty} \frac{2}{6x}.\]
    The second term goes to $0$. For the first term, apply L'H\^{o}pital's Rule:
    \[\lim_{x \to \infty} \frac{20\log x}{6x}
        \;=\; \lim_{x \to \infty} \frac{\tfrac{20}{x}}{6}\;=\; \lim_{x \to \infty} \frac{20}{6x}\;=\; 0.\]
    Therefore,
    \[\lim_{x \to \infty} \frac{20x\log x + 2x}{6x^2} \;=\; 0 + 0 \;=\; 0.\]
    Since the limit is $0 < \infty$, we conclude $f(x) \in O(g(x))$.
\end{solution}

\question[8] Answer the following parts:
\begin{parts}
\part Give pseudocode for an algorithm that finds the first repeated negative integer in a given sequence of integers.
\part Analyze the worst-case time complexity of the algorithm you devised in part (a). To receive full credit, you must state your choice of elementary operation.
\end{parts}

\begin{solution}\\
    \textbf{A:}
    \begin{verbatim}
   1 function FirstRepeatedNegative(numbers):
   2     seen = empty set
   3     for x in numbers:
   4     |    if x < 0:
   5     |        if x in seen:
   6     |            return x
   7     |        add x to seen
   8     end
   9     return "none"
    \end{verbatim}
    \textbf{B:}\\
    \textbf{Elementary operation:} Comparison (lines 3, 4, and 5).
    With $n$ elements this gives at most $3n$ comparisons.
    Worst-case time complexity: $O(n)$.
\end{solution}

\question[8] Answer the following parts:
\begin{parts}
\part Give pseudocode for an algorithm that determines whether a string of $n$
 characters has all unique characters.
\part Analyze the worst-case time complexity of the algorithm you devised in part (a). To receive full credit, you must state your choice of elementary operation.
\end{parts}

\begin{solution}\\
    \textbf{A:}
    \begin{verbatim}
   1 function UniqueString(s):
   2     seen = empty set
   3     for c in s:
   4     |    if c in seen:
   5     |        return false
   6     |    add c to seen
   7     end
   8     return true
    \end{verbatim}
    \textbf{B:}\\
    \textbf{Elementary operation:} Comparison (lines 3 and 4).
    With $n$ elements this gives at most $2n$ comparisons.
    Worst-case time complexity: $O(n)$.
\end{solution}

\question[8] Answer the following parts:
\begin{parts}
\part Give pseudocode for an algorithm that locates the last occurrence of the smallest element in a list of \(n\) integers. (Elements are not necessarily distinct.)
\part Analyze the worst-case time complexity of the algorithm you devised in part (a). To receive full credit, you must state your choice of elementary operation.
\end{parts}

\begin{solution}\\
    \textbf{A:}
    \begin{verbatim}
   1 function LastMinIndex(ints):
   2     min_val = ints[0]
   3     index = 0
   4     for i from 1 to n-1:
   5     |    if ints[i] <= min_val:
   6     |        min_val = ints[i]
   7     |        index = i
   8     end
   9     return index
    \end{verbatim}
    \textbf{B:}\\
    \textbf{Elementary operation:} Comparison (lines 4 and 5).
    The loop always runs from $i = 1$ to $i = n-1$ regardless of the input. Each iteration incurs 2 comparisons: the loop condition on line 4 and the $\leq$ check on line 5, giving $2(n-1)$ comparisons total.
    Worst-case time complexity: $O(n)$.
\end{solution}

\question[8] Answer the following parts:
\begin{parts}
\part Given an array of integers $nums$ of length $n$ and an integer target
, find three distinct indices $i,j,k$ such that $nums[i]+nums[j]+nums[k]=target$.

You may assume exactly one such triplet exists, and you may not use the same element more than once. Return the indices in any order.

\part Analyze the worst-case time complexity of the algorithm you devised in part (a). To receive full credit, you must state your choice of elementary operation.
\end{parts}

\begin{solution}\\
    \textbf{A:}
    \begin{verbatim}
   1 function ThreeSum(nums, target):
   2     for i from 0 to n-1:
   3     |    for j from i+1 to n-1:
   4     |    |    for k from j+1 to n-1:
   5     |    |    |    if nums[i] + nums[j] + nums[k] == target:
   6     |    |    |        return (i, j, k)
   7     |    |    end
   8     |    end
   9     end
    \end{verbatim}
 
    \textbf{B:}\\
    \textbf{Elementary operation:} Comparison (lines 2, 3, 4, and 5).
    The total number of comparisons is given by:
    \[\sum_{i=0}^{n-1}\sum_{j=i+1}^{n-1}\sum_{k=j+1}^{n-1} 4\]
    Evaluating the innermost sum:
    \[\sum_{k=j+1}^{n-1} 4 = 4(n-1-j)\]
    Substituting into the middle sum:
    \[\sum_{j=i+1}^{n-1} 4(n-1-j)= 4\sum_{m=0}^{n-2-i} m
        = 4\cdot\frac{(n-1-i)(n-2-i)}{2}= 2(n-1-i)(n-2-i)\]
    Substituting into the outer sum:
    \[\sum_{i=0}^{n-1} 2(n-1-i)(n-2-i)= 2\cdot\frac{n(n-1)(n-2)}{3}
        = \frac{2n(n-1)(n-2)}{3}= 4\cdot\binom{n}{3} \in \Theta(n^3)\]
    In the worst case every triple is examined, giving $\dfrac{2n(n-1)(n-2)}{3}$ total comparisons.
    Worst-case time complexity: $O(n^3)$.
\end{solution}
\end{questions}
\end{document}